% TARI29_Assignment_report_template.tex
% This is a template for the assignment reports
% TARI29 Artificial Intelligence, 7.5 credits, Winter 2024
% Original version by:  Vladimir Tarasov, 2024
% Revised by:           Alexandros Tzanetos, 2024 

% Use a modern class instead of an old article class
\documentclass{scrartcl}

\KOMAoptions{
    parskip=half,  % full, off
    fontsize=12pt, % base font size (10pt default)
    % headings=big,% small/normal/big headings (normal is default), 
    % paper=a5,    % paper format (a4 default) 
    pagesize=auto  % Use paper format for PDF too
}

% ---------------------- Report details --------------------- %
\newcommand{\docauthor}{Andrej Kutny, Charles Madjeri, Samrat Debnath \\ Vladimir Giustacchini, Aikeya Ainiwaer}
\newcommand{\courseyear}{2025}
\newcommand{\coursename}{TARI29 Artificial Intelligence}
\newcommand{\coursecredits}{7.5 credits}
\newcommand{\fullcoursename}{\coursename, \coursecredits}
\newcommand{\termname}{Autumn \courseyear}
\newcommand{\groupnumber}{4}
\newcommand{\reportname}{Group \groupnumber\ Assignment Report}
\newcommand{\assignmentname}{Evolutionary Computation Assignment}
% ------------------ End of report details ------------------ %



% ---------- Setting up properties of the PDF file ---------- %
\usepackage{hyperref}
\hypersetup{
    colorlinks=true,
    allcolors=blue,
    pdftitle={\reportname},
    pdfauthor={\docauthor},
    pdfsubject={\coursename, \termname}
	% pdfpagemode=FullScreen,
	% pdfpagemode= UseOutlines % the default if present
}
% ------------ End of properties of the PDF file ------------ %



% ----------- Setting up the headers and footers ------------ %
\usepackage{scrlayer-scrpage}
\pagestyle{scrheadings}
\KOMAoptions{% 
    headsepline,% line below the header
    plainheadsepline,% also on scrplain
}
\clearpairofpagestyles % Erase all current configurations (updated from deprecated \clearscrheadfoot)
% inner part of the header [titel_page]{other_pages}
\ihead[\coursename, \courseyear]{\coursename, \courseyear}
% outer part of the header [titel_page]{other_pages}
\ohead{\assignmentname}
% center part of the footer [titel_page]{other_pages}
\cfoot[\textup{Page \pagemark\ of \pageref{LastPage}}]{\textup{Page \pagemark\ of \pageref{LastPage}}}
% ------------- End of the headers and footers -------------- %



% ---------------- Setting up the title page ---------------- %
\title{\reportname}
\subtitle{An Evolutionary Algorithm for the Traveling Salesman Problem}
\author{\docauthor}
\date{\today}
% ------------------ End of the title page ------------------ %



% --------------- Packages and custom commands -------------- %
% Specify the input encoding of the source file as UTF8 
\usepackage[utf8]{inputenc}
% The T1 font encoding is an 8-bit encoding and fully supports words containing accented characters
\usepackage[T1]{fontenc}
% Load Latin modern fonts in outline format
\usepackage{lmodern}
% Load better typewriter font
\usepackage{inconsolata}
% For better hyphenation patterns
\usepackage[english]{babel}
% For improved micro-typography
\usepackage{microtype}
% Switch off extra space after a sentence
\frenchspacing
% To be able to iclude pictures
\usepackage{graphicx}
\usepackage{zref-totpages} % To get the total number of pages
\usepackage{lastpage} % For LastPage reference
\usepackage{mathtools,amsmath} % for advanced math formulas
 % To highlight source code
\usepackage{minted}
% Default options for all minted environments
\setminted{
           style=default, % for minted envronment
           autogobble, % automatically remove all common leading whitespace
           xleftmargin=10pt, % indentation to add (on the left) before the listing
           numbersep=6pt, % gap between numbers and start of line
           linenos, % enables line numbers
           mathescape % enables the usual math mode inside comments
}
\setmintedinline[c]{style=default} % for minted inlines
% Flexible handling of verbatim text
\usepackage{fancyvrb}
% For well-spaced lines and guidelines in tables
\usepackage{booktabs}
% To typeset algorithms or pseudocode in LaTeX 
\usepackage{algorithm}
\usepackage{algpseudocode}
% For plots and graphs
\usepackage{pgfplots}
\pgfplotsset{width=10cm,compat=1.18}

% To generate placeholder text - can be deleted when you have written real text
\usepackage{lipsum}
% ----------- End of packages and custom commands ----------- %


\begin{document}

\maketitle

% % It can be commented out if you do not want the table of contents
% \tableofcontents 


\section{Introduction}
\label{sec:intro}

\textcolor{red}{The introduction should briefly introduce the assignment and its purpose.}

\lipsum[4]


\section{The Traveling Salesman Problem}
\label{sec:problem_description}

\textcolor{red}{The second section should present the problem you tackle using your evolutionary approach. Overall, this section should include:}

{\color{red}
\begin{itemize}
    \item The mathematical formulation considered in your study. Some problems have a clear mathematical model (e.g., Travelling Salesman Problem), while others do not (e.g., $n$-Queens). Based on the problem you chose, search the literature and find a proper way to present the problem.
    \item One paragraph that briefly presents at least 3 published academic works where any evolutionary approach is used to solve the problem. It would be wise to cite here works that influenced your algorithm. This practice saves you time from looking for additional academic resources. You can find more information about reading and searching in the literature in \cite{zobel2014reading}.
    \item The motivation behind the evolutionary approach you decided to develop. A good practice would be to align the motivation with some literature gap found in the academic works you presented above. However, this is not mandatory. You can motivate your selection on the characteristics of the algorithm making it proper for the problem.
\end{itemize}
}

\lipsum[2]


\section{Algorithm}
\label{sec:algorithm}

\textcolor{red}{The third section should present the evolutionary approach you developed. You can divide this section into subsection. In any case, you should mention the following details:}

\textcolor{red}{\textbf{Evolutionary approach.} Clearly describe the algorithm you developed. You should clearly explain the evolutionary operators you used and what modifications you did to match the problem. It is extremely important to present also a pseudocode of your algorithm. An example is given in \ref{alg:pseudocode_example}, below. For more insight into presentation of algorithms, you can advise \cite{zobel2014algorithms}.}

{
\color{red}To typeset pseudocode in \LaTeX\ you can use one of the following options:
\begin{itemize}
    \item Choose ONE of the (\texttt{algpseudocode} OR \texttt{algcompatible} OR \texttt{algorithmic}) packages to typeset algorithm bodies, and the algorithm package for captioning the algorithm.
    \item The \texttt{algorithm2e} package.
\end{itemize}
You can find more information here: \url{https://www.overleaf.com/learn/latex/Algorithms}
}

\begin{algorithm}
\caption{Example of an algorithm's pseudocode}\label{alg:pseudocode_example}
\begin{algorithmic}
\Require $n \geq 0$
\Ensure $y = x^n$
\State $y \gets 1$
\State $X \gets x$
\State $N \gets n$
\While{$N \neq 0$}
\If{$N$ is even}
    \State $X \gets X \times X$
    \State $N \gets \frac{N}{2}$  \Comment{This is a comment}
\ElsIf{$N$ is odd}
    \State $y \gets y \times X$
    \State $N \gets N - 1$
\EndIf
\EndWhile
\end{algorithmic}
\end{algorithm}

\textcolor{red}{\textbf{Solution representation.} Clearly describe the solution representation you used. You can use figures to improve the comprehensibility of this part.}

\textcolor{red}{\textbf{Fitness function.} It is also very important to mention the fitness function you used. In many cases, the objective function of the problem is not the same as the fitness function used in an evolutionary algorithm. An example, following the principles of \cite{zobel2014mathematics}, is given below.}

\begin{equation}
    F = \sum_{i=1}^d x_i^2 
\end{equation}
where $x_i$ is the $i$-th gene (i.e., decision variable) in the solution and $d$ corresponds to the number of decision variables in the problem.

\textcolor{red}{\textbf{Note:} Change the section's title to match the name of the algorithm you developed for your assignment.}

\lipsum[3]


\section{Experimental part}
\label{sec:experimentation}

{\color{red}
This section describes the setup of experiments \cite{zobel2014experimentation}:

\begin{itemize}
    \item Provide the details of the hardware and software that you used.
    \item Describe the steps you carried out during your experiments.
    \item Detail the data you used for the evaluation of your algorithm.
\end{itemize}
}

This section describes the setup of experiments conducted to evaluate and compare the performance of the greedy best-first search baseline with two evolutionary algorithm variants employing different crossover operators.

\subsection{Hardware and Software Specifications}

\subsubsection{Hardware Configuration}
The experiments were conducted on the following hardware configurations:

\textbf{Configuration 1:}
\begin{itemize}
    \item \textbf{Processor}: AMD Ryzen 7 7840HS (16 cores) @ 3.8GHz
    \item \textbf{Memory}: 48 GB RAM
    \item \textbf{Operating System}: Linux Ubuntu 24.04
\end{itemize}

\textbf{Configuration 2:}
\begin{itemize}
    \item \textbf{Processor}: Apple M2 Pro (12 cores)
    \item \textbf{Memory}: 16 GB LPDDR5
    \item \textbf{Operating System}: macOS Tahoe 26.0.1
\end{itemize}

\subsubsection{Software Environment}
The algorithms were implemented using the following software tools:
\begin{itemize}
    \item \textbf{Programming Language}: Python 3.12
    \item \textbf{Libraries and Frameworks}: 
    \begin{itemize}
        \item NumPy (numerical computations and array operations)
        \item Matplotlib (data visualization and convergence plots)
    \end{itemize}
    \item \textbf{Development Environment}: Visual Studio Code
\end{itemize}

\subsection{Experimental Procedure}

\subsubsection{Algorithm Configuration}

\textbf{Design Choices:}
The evolutionary algorithm employs the following systematic design decisions:
\begin{itemize}
    \item \textbf{Selection Method}: Elite selection (fixed design choice ensuring the best individuals are preserved across generations)
\end{itemize}

\textbf{Configurable Parameters:}
The following parameters were varied during experimentation:
\begin{itemize}
    \item \textbf{Population Size}: \textcolor{red}{[TO BE COMPLETED - configurable parameter, e.g., 50, 100, 200]}
    \item \textbf{Stopping Criteria}: Configurable termination conditions (iteration-based, time-based, or improvement-based as detailed below)
    \item \textbf{Crossover Rate}: \textcolor{red}{[TO BE COMPLETED - configurable parameter, e.g., 0.7, 0.8, 0.9]}
    \item \textbf{Elite Rate}: \textcolor{red}{[TO BE COMPLETED - configurable parameter representing the proportion of population preserved as elite, e.g., 0.1, 0.2, 0.3]}
\end{itemize}

\subsubsection{Execution Protocol}
Each algorithm was executed following this protocol:
\begin{enumerate}
    \item \textbf{Independent Runs}: Each algorithm configuration was run at least 25 times with different random seeds to account for stochastic variation and ensure statistical validity. \textcolor{red}{[TO BE COMPLETED - specify if more runs were conducted based on stopping criteria]}
    \item \textbf{Initialization}: For evolutionary algorithms, the initial population was generated using random permutations of city sequences. For the greedy best-first search, different starting cities were used across runs to evaluate sensitivity to initial conditions.
    \item \textbf{Termination Criteria}: The algorithms employ multiple configurable stopping criteria:
    \begin{itemize}
        \item \textbf{Iteration-based}: Maximum number of generations/iterations \textcolor{red}{[TO BE COMPLETED - specify value when used]}
        \item \textbf{Time-based}: Maximum execution time limit \textcolor{red}{[TO BE COMPLETED - specify duration when used]}
        \item \textbf{Improvement-based}: Termination when no improvement in best solution is observed for a specified number of consecutive generations \textcolor{red}{[TO BE COMPLETED - specify threshold when used]}
    \end{itemize}
    The specific stopping criterion (or combination of criteria) used depends on the experimental configuration.
    \item \textbf{Performance Recording}: For each run, the best solution quality, average solution quality, and computational time were recorded.
\end{enumerate}

\subsection{Evaluation Data}

\subsubsection{Problem Instances}
The algorithms were evaluated on TSP instances of varying sizes to assess scalability and performance:
\begin{itemize}
    \item \textbf{Problem Sizes}: \textcolor{red}{[TO BE COMPLETED - e.g., 20, 50, 100, and 200 cities]}
    \item \textbf{Number of Instances}: \textcolor{red}{[TO BE COMPLETED - e.g., number of instances per problem size]}
\end{itemize}

\subsubsection{Data Generation}
The experimental evaluation employs two types of datasets:

\textbf{Benchmark Datasets}: Standard TSP instances from TSPLIB \textcolor{red}{[TO BE COMPLETED - specify which instances, e.g., berlin52, eil76, kroA100]}. These benchmark problems are parsed and stored as arrays of tuples containing city coordinates $(x, y)$.

\textbf{Self-Generated Datasets}: Randomly generated TSP instances \textcolor{red}{[TO BE COMPLETED - specify generation method, e.g., cities with coordinates uniformly distributed in a defined range, number of instances generated]}.

This combination of benchmark and synthetic datasets allows for both standardized comparison with existing literature and controlled experimentation across various problem characteristics.

\subsubsection{Distance Metric}
The distance metric used for calculating distances between cities is configurable as a parameter:

\textbf{Euclidean Distance}: 
\begin{equation}
d(i,j) = \sqrt{(x_i - x_j)^2 + (y_i - y_j)^2}
\end{equation}

\textbf{Manhattan Distance}: 
\begin{equation}
d(i,j) = |x_i - x_j| + |y_i - y_j|
\end{equation}

where $(x_i, y_i)$ and $(x_j, y_j)$ represent the coordinates of cities $i$ and $j$ respectively.

\subsection{Performance Metrics}

The following metrics were used to evaluate and compare algorithm performance:

\begin{enumerate}
    \item \textbf{Best Solution Quality}: The length of the shortest tour found across all runs for each algorithm.
    
    \item \textbf{Average Solution Quality}: The mean tour length across all independent runs, providing insight into algorithm consistency.
    
    \item \textbf{Standard Deviation}: The variability of solution quality across runs, indicating algorithm robustness.
    
    \item \textbf{Computational Time}: The average execution time required for each algorithm to complete.
    
    \item \textbf{Convergence Behavior}: The evolution of solution quality over generations/iterations, analyzed through convergence plots.
\end{enumerate}

Statistical significance of performance differences between algorithms was assessed using \textcolor{red}{[TO BE COMPLETED - e.g., paired t-tests with $\alpha = 0.05$ significance level]}.


\section{Results and Analysis}
\label{sec:results-analysis}

\textcolor{red}{This section should present the obtained results and provide an insightful analysis of them. You can present the results using graphs, tables, or any other visualization method suits your purpose. Do not forget to include proper captions \cite{zobel2014graphs} in any of these illustration methods you use. You do not need to provide any execution details as they are already presented in Sec.~\ref{sec:experimentation}.}

\textcolor{red}{A good practise would be to compare your algorithm with a simpler approach, such as (a) a naive method, (b) a Hill Climbing approach, or (c) a simple evolutionary algorithm. In the third case, you can use the simpler version of the algorithm you developed, i.e., the original algorithm without your modifications. In that case, you should briefly describe the comparing method(s) in Sec.~\ref{sec:experimentation}. Alternatively, you can use some reference results derived from the repositories you found some benchmark instances.}

\textcolor{red}{To display tables, the \texttt{booktabs} package might be useful. For example, Table~\ref{tab:results_example} shows how you should increase the  size of $n$, when running your code. You can advice \cite{zobel2014graphs} to see a few examples of proper tables.}

\begin{table}[ht]
\centering
\caption{Example of comparison the developed algorithm's results with the best ones from a repository.}
\label{tab:results_example}
	\begin{tabular}{lrrr}
	\toprule
	\textbf{Instance} & \textbf{Optimum (Repository xyz)} & \textbf{EA} & \textbf{time (s)}\\
	\midrule
	st70       &  678.597      & 677.109	& 0.67\\
	ei176      &  545.387      & 544.369	& 1.16\\
	kroA100    &  21285.443    & 21285.443	& 1.69\\
	rd100      &  7910.396     & 7910.396	& 2.14\\
	Pr136      &  96772        & 96770.924	& 7.11\\
	Pr144      &  58537        & 58535.221	& 7.97\\
	a280       &  2856.769     & 2856.769	& 33.47\\
	\bottomrule
	\end{tabular}
\end{table}

\textcolor{red}{You can use different illustration methods to present different aspects of your analysis. Figure~\ref{fig:plot_example} gives an example using the \href{https://www.overleaf.com/learn/latex/Pgfplots_package}{\texttt{pgfplots}} package.}

\begin{figure}[ht]
\centering
\begin{tikzpicture}
\begin{axis}[
    % title={Example of convergence analysis},
    xlabel={Generation},
    ylabel={Fitness function value},
    xmin=0, xmax=20,
    ymin=290, ymax=450,
    xtick={0,5,10,15,20},
    ytick={290,300,350,400,450},
    legend pos=north east,
    ymajorgrids=true,
    grid style=dashed,
]

\addplot[
    color=blue,
    mark=square,
    ]
    coordinates {
    (0,420)(4,411)(7,387)(10,382)(14,364)(18,360)(20,358)
    };
    \addlegendentry{EA1}

\addplot[
    color=red,
    mark=triangle,
    ]
    coordinates {
    (0,422)(4,373)(9,362)(12,312)(18,311)(20,309)
    };
    \addlegendentry{EA2}

\addplot[
    color=violet,
    mark=diamond,
    ]
    coordinates {
    (0,448)(6,401)(8,349)(11,325)(14,299)(20,298)
    };
    \addlegendentry{EA3}
  
\end{axis}
\end{tikzpicture}
\caption{Example of convergence analysis.}
\label{fig:plot_example}
\end{figure}

\lipsum[4]

\section{Conclusions}
\label{sec:conclusions}

\textcolor{red}{In this section you should provide a concise summary of what has been done, the obtained results and some recommendations on how this study could be extended.}

\lipsum[4]

\bibliographystyle{ieeetr}
\bibliography{references}

\end{document}
